\documentclass[10pt,a4paper]{amsart}
\usepackage{amsmath,amssymb,amsthm,graphicx}
% Retained preamble for the Stacks Project chapter snapshot in this directory.
% The chapter's own preamble.tex is a copy of the shared Stacks Project preamble; the block
% below is sliced MECHANICALLY, line for line, from it. Nothing was added, removed or
% redefined; the wrapper class, the xr-hyper/hyperref wiring and the externaldocument list
% are omitted because the assembler supplies the document class.
%
% lines taken from preamble.tex: 3, 5, 9, 11, 12, 156, 157, 158, ... (51 lines total)
%
% package coverage restored from the chapter's own preamble (2026-10-13): xy, xy-2cell

\usepackage[all]{xy}
\xyoption{2cell}
\UseAllTwocells

\newenvironment{reference}{\comment}{\endcomment}
\newenvironment{slogan}{\comment}{\endcomment}
\newenvironment{history}{\comment}{\endcomment}
\theoremstyle{plain}
\newtheorem{theorem}[subsection]{Theorem}
\newtheorem{proposition}[subsection]{Proposition}
\newtheorem{lemma}[subsection]{Lemma}
\theoremstyle{definition}
\newtheorem{definition}[subsection]{Definition}
\newtheorem{example}[subsection]{Example}
\newtheorem{exercise}[subsection]{Exercise}
\newtheorem{situation}[subsection]{Situation}
\theoremstyle{remark}
\newtheorem{remark}[subsection]{Remark}
\newtheorem{remarks}[subsection]{Remarks}
\numberwithin{equation}{subsection}
\def\lim{\mathop{\mathrm{lim}}\nolimits}
\def\colim{\mathop{\mathrm{colim}}\nolimits}
\def\Spec{\mathop{\mathrm{Spec}}}
\def\Hom{\mathop{\mathrm{Hom}}\nolimits}
\def\Ext{\mathop{\mathrm{Ext}}\nolimits}
\def\SheafHom{\mathop{\mathcal{H}\!\mathit{om}}\nolimits}
\def\SheafExt{\mathop{\mathcal{E}\!\mathit{xt}}\nolimits}
\def\Sch{\mathit{Sch}}
\def\Mor{\mathop{\mathrm{Mor}}\nolimits}
\def\Ob{\mathop{\mathrm{Ob}}\nolimits}
\def\Sh{\mathop{\mathit{Sh}}\nolimits}
\def\NL{\mathop{N\!L}\nolimits}
\def\CH{\mathop{\mathrm{CH}}\nolimits}
\def\proetale{{pro\text{-}\acute{e}tale}}
\def\etale{{\acute{e}tale}}
\def\QCoh{\mathit{QCoh}}
\def\Ker{\mathop{\mathrm{Ker}}}
\def\Im{\mathop{\mathrm{Im}}}
\def\Coker{\mathop{\mathrm{Coker}}}
\def\Coim{\mathop{\mathrm{Coim}}}
\DeclareMathSymbol{\boxtimes}{\mathbin}{AMSa}{"02}
\def\QCohstack{\mathcal{QC}\!\mathit{oh}}
\def\Cohstack{\mathcal{C}\!\mathit{oh}}
\def\Spacesstack{\mathcal{S}\!\mathit{paces}}
\def\Quotfunctor{\mathrm{Quot}}
\def\Hilbfunctor{\mathrm{Hilb}}
\def\Curvesstack{\mathcal{C}\!\mathit{urves}}
\def\Polarizedstack{\mathcal{P}\!\mathit{olarized}}
\def\Complexesstack{\mathcal{C}\!\mathit{omplexes}}
\def\Pic{\mathop{\mathrm{Pic}}\nolimits}
\def\Picardstack{\mathcal{P}\!\mathit{ic}}
\def\Picardfunctor{\mathrm{Pic}}
\def\Deformationcategory{\mathcal{D}\!\mathit{ef}}
\usepackage{cleveref}
\crefname{originallemma}{original lemma}{original lemmas}
\Crefname{originallemma}{Original lemma}{Original lemmas}
\makeatletter
\newlabel{lemma-algebraization-principal-bis}{{0EJ9}{source}{}{sourcepointer-1}{}}
\newlabel{lemma-algebraization-principal-bis@cref}{{[originallemma][0][]0EJ9}{[1][0][]source}{}{}{}}
\makeatother
\crefname{originallemma}{original lemma}{original lemmas}
\Crefname{originallemma}{Original lemma}{Original lemmas}
\makeatletter
\newlabel{lemma-fully-faithful-general}{{0EKU}{source}{}{sourcepointer-2}{}}
\newlabel{lemma-fully-faithful-general@cref}{{[originallemma][0][]0EKU}{[1][0][]source}{}{}{}}
\makeatother
\crefname{originallemma}{original lemma}{original lemmas}
\Crefname{originallemma}{Original lemma}{Original lemmas}
\makeatletter
\newlabel{lemma-unwinding-conditions-bis}{{0EJB}{source}{}{sourcepointer-3}{}}
\newlabel{lemma-unwinding-conditions-bis@cref}{{[originallemma][0][]0EJB}{[1][0][]source}{}{}{}}
\makeatother
\crefname{originallemma}{original lemma}{original lemmas}
\Crefname{originallemma}{Original lemma}{Original lemmas}
\makeatletter
\newlabel{proposition-cd-1}{{0EJJ}{source}{}{sourcepointer-4}{}}
\newlabel{proposition-cd-1@cref}{{[originallemma][0][]0EJJ}{[1][0][]source}{}{}{}}
\makeatother
\crefname{originallemma}{original lemma}{original lemmas}
\Crefname{originallemma}{Original lemma}{Original lemmas}
\makeatletter
\newlabel{situation-algebraize}{{0EHC}{source}{}{sourcepointer-5}{}}
\newlabel{situation-algebraize@cref}{{[originallemma][0][]0EHC}{[1][0][]source}{}{}{}}
\makeatother
\title{Under dualizing-complex, principal-ideal and depth-dimension hypotheses the completion functor on finite locally free sheaves is an equivalence}
\author{The Stacks Project Authors}
\date{stacks-project:algebraization.tex:local-snapshot-20261012; GNU FDL 1.2 or later}
\begin{document}
\maketitle
\noindent\textit{Editorial note.} Statement span is lines 8089-8113 of sources/stacks-algebraization/pinned-original-source.tex; proof span is lines 8115-8131, the complete author proof (15 body lines).\par\medskip
\noindent\textit{Editorial note.} Navigation stubs: 5 cited-but-not-included labels are supplied as non-load-bearing pointers with their pinned tags where the snapshot tag index has them.\par\medskip
\noindent\textit{Editorial note.} extra\_packages is empty on purpose: the retained preamble already loads xy with options and a second declaration clashes in the assembler wrapper.\par\medskip
\noindent\textit{Editorial note.} Text is retained verbatim, including the author's own typos and placeholders.\par\medskip
\noindent\textit{Editorial note.} Pool algebraization wave36 (al01 batch). Spans per computed label truth (al\_clean\_queue), not pool-card lines.\par\medskip
\noindent\textit{Editorial note.} Cross-chapter homonym warning recorded in dedup (chapter-aware gate).\par\medskip
\noindent\textit{Editorial note.} Cross-chapter same-label warning (NOT an occupancy conflict, Chief ruling 2026-10-13 evening): label 'lemma-equivalence-better' also exists in stacks-pione — a different proposition in a different chapter file.\par\medskip
\noindent\textit{Allowed standard prerequisites (editorial).}\begin{itemize}
\item Explicit prerequisite boundary. See why-field; in-chapter lemmas cited by name inside the retained author proof become stub prereqs.
\end{itemize}\medskip
\ifdefined\hypertarget\hypertarget{sourcepointer-1}{}\fi
\noindent\textit{Source reference.} Original lemma 0EJ9: sources/stacks-algebraization/pinned-original-source.tex, source lines 6736--6744. Cited by the retained statement or proof but not part of the collected spans; supplied only as a navigation pointer. The argument does not depend on its content.\par\medskip
\ifdefined\hypertarget\hypertarget{sourcepointer-2}{}\fi
\noindent\textit{Source reference.} Original lemma 0EKU: sources/stacks-algebraization/pinned-original-source.tex, source lines 5062--5075. Cited by the retained statement or proof but not part of the collected spans; supplied only as a navigation pointer. The argument does not depend on its content.\par\medskip
\ifdefined\hypertarget\hypertarget{sourcepointer-3}{}\fi
\noindent\textit{Source reference.} Original lemma 0EJB: sources/stacks-algebraization/pinned-original-source.tex, source lines 6939--6957. Cited by the retained statement or proof but not part of the collected spans; supplied only as a navigation pointer. The argument does not depend on its content.\par\medskip
\ifdefined\hypertarget\hypertarget{sourcepointer-4}{}\fi
\noindent\textit{Source reference.} Original lemma 0EJJ: sources/stacks-algebraization/pinned-original-source.tex, source lines 7480--7483. Cited by the retained statement or proof but not part of the collected spans; supplied only as a navigation pointer. The argument does not depend on its content.\par\medskip
\ifdefined\hypertarget\hypertarget{sourcepointer-5}{}\fi
\noindent\textit{Source reference.} Original lemma 0EHC: sources/stacks-algebraization/pinned-original-source.tex, source lines 5200--5204. Cited by the retained statement or proof but not part of the collected spans; supplied only as a navigation pointer. The argument does not depend on its content.\par\medskip
\begin{lemma}
\label{lemma-equivalence-better}
In Situation \ref{situation-algebraize} assume
\begin{enumerate}
\item $A$ has a dualizing complex and is $I$-adically complete,
\item $I = (f)$ generated by a single element,
\item $A$ is local with maximal ideal $\mathfrak a = \mathfrak m$,
\item one of the following is true
\begin{enumerate}
\item $A_f$ is $(S_2)$ and for $\mathfrak p \subset A$,
$f \not \in \mathfrak p$ minimal we have $\dim(A/\mathfrak p) \geq 4$, or
\item if $\mathfrak p \not \in V(f)$ and
$V(\mathfrak p) \cap V(f) \not = \{\mathfrak m\}$, then
$\text{depth}(A_\mathfrak p) + \dim(A/\mathfrak p) > 3$.
\end{enumerate}
\end{enumerate}
Then with $U_0 = U \cap V(f)$ the completion functor
$$
\colim_{U_0 \subset U' \subset U\text{ open}}
\textit{Coh}(\mathcal{O}_{U'})
\longrightarrow
\textit{Coh}(U, f\mathcal{O}_U)
$$
is an equivalence on the full subcategories of finite locally free objects.
\end{lemma}

\begin{proof}
It follows from Lemma \ref{lemma-fully-faithful-general}
that the functor is fully faithful (details omitted).
Let us prove essential surjectivity. Let $(\mathcal{F}_n)$ be a finite locally
free object of $\textit{Coh}(U, f\mathcal{O}_U)$. By either
Lemma \ref{lemma-algebraization-principal-bis} or
Proposition \ref{proposition-cd-1}
there exists a coherent $\mathcal{O}_U$-module $\mathcal{F}$
such that $(\mathcal{F}_n)$ is the completion of $\mathcal{F}$.
Namely, for the application of either result the only thing to
check is that $(\mathcal{F}_n)$ satisfies the $(2, 3)$-inequalities.
This is done in Lemma \ref{lemma-unwinding-conditions-bis}. If $y \in U_0$,
then the $f$-adic completion of the stalk $\mathcal{F}_y$ is isomorphic to
a finite free module over the $f$-adic completion of $\mathcal{O}_{U, y}$.
Hence $\mathcal{F}$ is finite locally free in an open neighbourhood
$U'$ of $U_0$. This finishes the proof.
\end{proof}

\end{document}
